Um dos primeiros passos na exploração de uma base de dados é a identificação do tipo dos
dados que estamos lidando a fim de empregar as técnicas exploratórias mais adequadas a
cada caso.
Os dados podem ser vir nos mais variados fomatos, como métricas, categoriais, nomes,
datas, etc.
No entanto, por mais variados que sejam os formatos possíveis, normalmente classicamos-os
de acordo com um número limitado de tipos.
Após isso, selecionamos quais técnicas usaremos para representar, apresentar e visualizar
os dados visando a identificação de padrões e a geração das primeiras hipóteses que irão
guiar o resto da exploração.
